\subsection{Analyser une équipe}

\begin{UPSTIManipulation}{Étape 9 — Qui frappe le plus fort ?}
  \textbf{Contexte :} Avant un combat, on veut désigner comme attaquant principal le personnage
  de l'équipe ayant l'attaque la plus élevée.

  \textbf{Cahier des charges :}
  \begin{itemize}
    \item Écrire une fonction qui reçoit l'équipe (tableau de structures et sa taille) et
      renvoie l'\textbf{indice} du personnage ayant l'attaque la plus élevée (en cas d'égalité,
      garder le premier trouvé)
    \item Afficher l'indice et le nom du personnage le plus offensif
  \end{itemize}

  \textbf{Exemple d'exécution :}
  \begin{lstlisting}[language=bash]
Arthor (PV: 35, ATQ: 7, DEF: 6)
Lysia (PV: 20, ATQ: 10, DEF: 2)
Finn (PV: 25, ATQ: 9, DEF: 3)
Attaquant principal: indice 1 (Lysia, ATQ=10)
  \end{lstlisting}

  \UPSTIetape{Modifier \texttt{jeu\_combat.c} pour respecter le cahier des charges ci-dessus.}
\end{UPSTIManipulation}

\subsection{Un butin après le combat}

\begin{UPSTIManipulation}{Étape 10 — Ramasser un objet}
  \textbf{Contexte :} Après un combat, les monstres vaincus laissent tomber des objets,
  représentés par un tableau de structures \texttt{Objet} (indépendant de
  \texttt{Personnage}), parmi lesquels on cherche le meilleur à équiper.

  \textbf{Cahier des charges :}
  \begin{itemize}
    \item Déclarer une structure \texttt{Objet} avec un nom et un bonus d'attaque
    \item Constituer un tableau de 3 objets (saisie du nom et du bonus pour chacun)
    \item Trouver l'objet ayant le plus grand bonus d'attaque, et équiper un personnage avec
      (augmenter son attaque du bonus de l'objet)
    \item Afficher la fiche du personnage avant et après l'équipement, ainsi que le nom de
      l'objet équipé
  \end{itemize}

  \textbf{Exemple d'exécution :}
  \begin{lstlisting}[language=bash]
Objet 0 - Nom: Epee rouillee
Objet 0 - Bonus d'attaque: 1
Objet 1 - Nom: Lame des braves
Objet 1 - Bonus d'attaque: 5
Objet 2 - Nom: Dague usee
Objet 2 - Bonus d'attaque: 2

Avant : Arthor (PV: 35, ATQ: 7, DEF: 6)
Objet equipe : Lame des braves (bonus +5)
Apres : Arthor (PV: 35, ATQ: 12, DEF: 6)
  \end{lstlisting}

  \UPSTIetape{Modifier \texttt{jeu\_combat.c} pour respecter le cahier des charges ci-dessus.}
\end{UPSTIManipulation}

\subsection{Se soigner}

\begin{UPSTIManipulation}{Étape 11 — Boire une potion}
  \textbf{Contexte :} On introduit des potions qui produisent des effets différents : soin ou
  poison. Le joueur doit pouvoir \textbf{décider}, à son tour, d'utiliser une potion plutôt que
  d'attaquer : ce n'est plus un déclenchement automatique.

  \textbf{Cahier des charges :}
  \begin{itemize}
    \item Représenter un type d'effet (soin ou poison) et une structure \texttt{Potion} (nom,
      effet, valeur)
    \item Écrire une fonction qui applique l'effet d'une potion à un personnage :
      \begin{itemize}
        \item soin : augmente les PV, sans dépasser 50
        \item poison : diminue les PV, sans descendre sous 0
      \end{itemize}
    \item Faire subir des dégâts à un personnage, puis proposer au joueur un menu à deux choix
      (par exemple \texttt{1. Attaquer}, \texttt{2. Utiliser une potion}), lu au clavier
    \item Selon le choix du joueur, soit le personnage attaque (comme à l'étape 6), soit il
      boit une potion de soin ; afficher sa fiche avant et après l'action choisie
  \end{itemize}

  \textbf{Exemple d'exécution :}
  \begin{lstlisting}[language=bash]
Avant l'action : Arthor (PV: 10, ATQ: 7, DEF: 6)
1. Attaquer
2. Utiliser une potion
Votre choix: 2
Arthor boit une Potion de soin (+20)
Apres l'action : Arthor (PV: 30, ATQ: 7, DEF: 6)
  \end{lstlisting}

  \UPSTIetape{Modifier \texttt{jeu\_combat.c} pour respecter le cahier des charges ci-dessus.}
\end{UPSTIManipulation}

\subsection{Classer son équipe}

\begin{UPSTIManipulation}{Étape 12 — Classement par dégâts infligés}
  \textbf{Contexte :} À la fin d'un combat, on veut classer les membres de l'équipe du plus au
  moins offensif, selon les dégâts qu'ils ont infligés.

  \textbf{Cahier des charges :}
  \begin{itemize}
    \item Ajouter à chaque personnage un total de dégâts infligés (saisi directement, sans
      simuler tout le combat)
    \item Trier l'équipe par ordre décroissant de dégâts infligés, en adaptant le tri par
      sélection vu à la séquence précédente
    \item Afficher le classement obtenu
  \end{itemize}

  \textbf{Exemple d'exécution :}
  \begin{lstlisting}[language=bash]
Personnage 0 - Nom: Arthor
Personnage 0 - Degats infliges: 42
Personnage 1 - Nom: Lysia
Personnage 1 - Degats infliges: 65
Personnage 2 - Nom: Finn
Personnage 2 - Degats infliges: 18

Classement :
1. Lysia : 65
2. Arthor : 42
3. Finn : 18
  \end{lstlisting}

  \UPSTIetape{Modifier \texttt{jeu\_combat.c} pour respecter le cahier des charges ci-dessus.}
\end{UPSTIManipulation}
